\addtocontents{toc}{\protect\Needspace{3\baselineskip}}
\originalchapter{守恒流、Regge 轨迹与共形参数}\label{ch:basic-charges}
对应 8.251 第 13 讲后半、第 14--15 讲. 守恒方程是局部结论, 守恒荷还要检查边界通量. 然后用弦的能量与角动量定义 $\alpha'$, 并明确改用无量纲共形参数的归一化.

\originalsection{Noether 定理的一次积分分部}
对场 $\phi^A(\xi)$ 的拉格朗日密度 $\mathcal L(\phi,\partial\phi)$, 设内部连续变换 $\delta\phi^A=\epsilon h^A$ 满足 $\delta\mathcal L=\epsilon\partial_aB^a$. 记 $\Pi_A^a=\partial\mathcal L/\partial(\partial_a\phi^A)$. 链式法则与欧拉--拉格朗日方程给
\[
 \frac{\delta\mathcal L}{\epsilon}
 =\frac{\partial\mathcal L}{\partial\phi^A}h^A+\Pi_A^a\partial_ah^A
 =\partial_a(\Pi_A^ah^A).
\]
所以 $J^a=\Pi_A^ah^A-B^a$ 满足 $\partial_aJ^a=0$. 原课主要用严格不变 $B^a=0$ 的情况. 对 $Q=\int J^0\dd^p\xi$, 散度定理给 $\dot Q=-\int_{\partial\Sigma}\boldsymbol J\cdot\dd\boldsymbol a$. 只有边界通量为零或配对抵消时才有守恒荷.

在粒子力学中, 同样证明 $Q=p_Ah^A-B$ 守恒. 例如 $L$ 不依 $q$ 时, $h=1,B=0$, 所得 $Q=p$; 这说明平移对称与动量守恒是同一个结构. 电磁连续方程 $\partial_t\rho+\nabla\cdot\boldsymbol J=0$ 则是局部电荷版.

\originalsection{弦的目标平移与截面独立性}
对 $\delta X^\mu=a^\mu$ 的常数平移, 守恒流正是 $P_\mu^a$. 更一般的世界面空间截线 $\gamma$ 上定义
\[
 p_\mu(\gamma)=\int_\gamma(P_\mu^\tau\dd\sigma-P_\mu^\sigma\dd\tau).
\]
该一形式的外微分为 $(\partial_\tau P_\mu^\tau+\partial_\sigma P_\mu^\sigma)\dd\tau\wedge\dd\sigma=0$. 对两条截线之间的区域用 Stokes 定理, 加上边界无通量, 得 $p_\mu(\gamma_1)=p_\mu(\gamma_2)$. 因此荷不依赖选哪个良好的世界面时间切片, 不是只在 $\tau$ 恰好等于目标时间时才守恒. 闭弦切线还应取同一同调类, 不把绕世界面不同次数的曲线当同一截面.

\originalsection{洛伦兹流与角动量}
无穷小洛伦兹变换为 $\delta X^\mu=\omega^\mu{}_{\nu}X^\nu$, $\omega_{\mu\nu}=-\omega_{\nu\mu}$. 对任意两个切向量, 内积的一阶变化为反对称 $\omega$ 与对称乘积缩并, 故 $\delta h_{ab}=0$. 由 Noether 定理得到
\[
 m^{a\mu\nu}=X^\mu P^{a\nu}-X^\nu P^{a\mu},\qquad
 M^{\mu\nu}=\int m^{\tau\mu\nu}\dd\sigma.
\]
也可直接求散度: 含 $\partial_aP$ 的两项用方程消失, 剩下 $\partial_aX^\mu P^{a\nu}-(\mu\leftrightarrow\nu)$, 因 $P$ 由对称 $h^{ab}$ 与切向量线性组合而为零. $M^{ij}$ 是通常角动量, $M^{0i}$ 是 boost 荷. 固定 D 膜背景只保留与边界条件相容的洛伦兹变换, 其余荷可向膜流失.

\originalsection{旋转直弦的 Regge 关系}
使用第\ref{ch:basic-motion}章的旋转解, 有
\[
 J=M^{12}=\int_0^{\pi a}\dd\sigma\,
 \frac{T_0a}{c}\cos^2(\sigma/a)
 =\frac{\pi T_0a^2}{2c}=\frac{E^2}{2\pi T_0c}.
\]
这是经典最大伸长旋转弦的角动量与能量关系. $J/\hbar$ 为无量纲, 因而若用能量单位定义斜率 $\alpha'_E$, 则
\[
 J/\hbar=\alpha'_E E^2,\qquad
 \alpha'_E=\frac1{2\pi T_0c\hbar}.
\]
用长度平方定义通常的弦参数为
\[
 \alpha'= (\hbar c)^2\alpha'_E
 =\frac{\hbar c}{2\pi T_0},\qquad\ell_s=\sqrt{\alpha'}.
\]
进入自然单位后这些记法归并, $T_s=1/(2\pi\alpha')$, $J=\alpha'E^2$. 量子理论会加入截距和选择规则, 不能把这条经典关系直接当所有弦态的精确质量公式.

\originalsection{统一的无量纲共形规范}
从此取 $c=\hbar=1$, 并重选无量纲 $\tau,\sigma$. 用常协向量 $n_\mu$ 选演化方向, 在 $n\cdot p\ne0$ 的区域定
\[
 n\cdot X=\beta\alpha'(n\cdot p)\tau,
 \quad 0\leq\sigma\leq\frac{2\pi}{\beta},
 \quad n\cdot P^\tau=\frac{\beta(n\cdot p)}{2\pi},
 \quad n\cdot P^\sigma=0,
\]
其中开弦 $\beta=2$, 闭弦 $\beta=1$. 第一式给 $n\cdot X'=0$. 第四式结合 Nambu--Goto 动量流给 $\dot X\cdot X'=0$. 将第二个密度条件代入 $P^\tau$ 后, 在非退化内点得到
\[
 \dot X^2+X'^2=0,\qquad
 P^{\tau\mu}=\frac{\dot X^\mu}{2\pi\alpha'},
 \qquad P^{\sigma\mu}=-\frac{X'^\mu}{2\pi\alpha'}.
\]
于是 $\ddot X-X''=0$, 同时必须保留 $(\dot X\pm X')^2=0$. 此过程固定了参数的尺度, 前一章有长度单位的 $\sigma$ 不再是这里的 $\sigma$. 开弦的范围 $\pi$ 与闭弦的 $2\pi$ 正是振子归一化不同的来源.

\originalsection{选择时间与光锥方向}
若 $n$ 选取普通时间分量, 上述规范与能量参数版本只差常数缩放. 若 $n\cdot X=X^+$, 就得到开弦光锥规范 $X^+=2\alpha'p^+\tau$ 和闭弦 $X^+=\alpha'p^+\tau$. 可把初始常数设零以选时间原点. 横向 $X^I$ 是待求的独立坐标, $X^-$ 将由约束决定.

对开弦, 零 $n\cdot P^\sigma$ 既来自端点无通量, 也可由局部方程与均匀密度延拓到内部. 对闭弦没有端点条件可借用, 必须额外选择参数标记线使 $\dot X\cdot X'=0$. 闭弦仍留有整体平移 $\sigma\mapsto\sigma+\sigma_0$ 的自由度, 它与下一章的全局周期约束相关.

我们在本编通过 Nambu--Goto 直接整理共形参数; 第二编将引入独立世界面度规的 Polyakov 形式, 更适于讨论路径积分和 Weyl 对称性. 两者在经典非退化条件下等价, 量子异常问题需要另行检验.
